% Part: proof-theory
% Chapter: cut-elimination
% Section: ce-topmost

\documentclass[../../../include/open-logic-section]{subfiles}

\begin{document}

\olfileid{pt}{cut}{top}

\olsection{移除最顶层的切割}

\begin{defn}
考虑一个以 \CutCS{} 规则结束的推导~$\pi$，
\[
\AxiomC{}
\RightLabel{$\pi_1$}
\Deduce$\Gamma \fCenter \Delta, !A$
\AxiomC{}
\RightLabel{$\pi_2$}
\Deduce$!A, \Gamma \fCenter \Delta$
\RightLabel{\CutCS}
\BinaryInf$\Gamma \fCenter \Delta$
\DisplayProof
\]
且不含其他 \CutCS{} 推理。$\pi$ 的\emph{切割高度}
为 $\cheight{\pi} = \pheight{\pi_1} + \pheight{\pi_2}$。$\pi$ 的
\emph{切割秩}为 $\cutrank{\pi} = \depth{!A}$。
\end{defn}

\begin{lem}\ollabel{lem:cut-adm-G3c}
\CutCS{} 规则在~\Log{G3c} 中是可接纳的。
换言之，若推导~$\pi$ 以单一~\CutCS{} 规则结束，且其余部分仅使用~\Log{G3c} 的规则，
则在~\Log{G3c} 中存在具有相同末相继式的推导~$\pi'$。
\end{lem}

\begin{proof}
证明使用对切割高度和秩的双重归纳。推导~$\pi$ 以如下推理结束：
\[
\AxiomC{}
\RightLabel{$\pi_1$}
\Deduce$\Gamma \fCenter \Delta, !A$
\AxiomC{}
\RightLabel{$\pi_2$}
\Deduce$!A, \Gamma \fCenter \Delta$
\RightLabel{\CutCS}
\BinaryInf$\Gamma \fCenter \Delta$
\DisplayProof
\]

归纳基础：$\cheight{\pi} = 0$ 且 $\cutrank{\pi} = 0$。由
$\cheight{\pi} = \pheight{\pi_1} + \pheight{\pi_2} = 0$可知，$\Gamma
\Sequent \Delta, !A$ 和 $!A, \Gamma \Sequent \Delta$ 都是公理。
有两种情况： 或(b)~$!A$ 在两个公理中都是主公式。在两种情况下，下面都将证明 $\Gamma \Sequent \Delta$ 必为公理（且无需使用归纳假设）。

现在，我们根据 $\pi_1$ 和 $\pi_2$ 是公理还是以某个推理结束，以及 $!A$~是否是该推理的主公式，来区分多种可能的情况。情况 A 和~B 涵盖了归纳基础。在归纳步骤中，我们假设 $\cheight{\pi} > 0$
或 $\cutrank{\pi} > 0$（或两者都成立）。我们可以使用归纳假设，即：对于任何以单个~\CutCS{} 结尾的推导，若该~\CutCS{} 的切割秩等于 $\cutrank{\pi}$ 且切割高度小于 $\cheight{\pi}$，或其切割秩小于 $\cutrank{\pi}$，则都可以从中移除最后的~\CutCS{}。$\cheight{\pi} = 0$ 的情况（两个前提都是公理）已由情况~A 和~B 涵盖。$\cheight{\pi} >
0$ 的情况则由情况 C--D 涵盖。

\paragraph{情况 A：}
某个 $!B$ 同时出现在 $\Gamma$ 和~$\Delta$ 中。
那么 $\Gamma \Sequent \Delta$ 是公理（如果 $!B$~是原子式），或者根据
\olref[seq][ex]{prop:G3c-axioms}，在 \Log{G3c} 中存在一个不含 \CutCS{} 的推导~$\pi'$。这涵盖了归纳基础的情况~(a)，因为如果 $!A$ 在 $\Gamma
\Sequent \Delta, !A$ 中不是主公式，或在 $!A, \Gamma \Sequent \Delta$ 中不是主公式，并且这些相继式是公理，那么必定有另一个原子式~$!B$ 同时出现在 $\Gamma$ 和~$\Delta$ 中。在归纳步骤中，如果有一个前提是公理且其中 $!A$ 不是主公式（即使另一个前提是某个规则的结论），该情况也适用。

\paragraph{情况 B：} 两个前提都是公理，且 $!A$ 是主公式，因而是原子式。考虑左前提 $\Gamma \Sequent
\Delta, !A$。由于 $!A$ 是主公式，它必然出现在 $\Gamma$ 中（否则 $\Gamma \Sequent \Delta, !A$ 就不是公理）。类似地，$!A$ 必然出现在 $\Delta$ 中，否则 $!A, \Gamma \Sequent \Delta$
就不是公理。因此 $!A$ 是原子式且同时出现在 $\Gamma$ 和 $\Delta$ 中，即 $\Gamma \Sequent \Delta$ 是公理。这涵盖了归纳基础的情况~(b)。

\begin{center}
替代的情况 A 与 B
\end{center}

\paragraph{情况 A：} 切割的某个前提是形如 $!C, \Gamma' \Sequent \Delta', !C$ 的公理，其中 $!C$~是原子式。如果 $!C$ 不是切割公式~$!A$，那么 $!C$ 必然同时出现在 $\Gamma$
和~$\Delta$ 中。此时，$\Gamma \Sequent \Delta$ 是公理，我们可以直接取它为~$\pi'$。如果 $!C \ident !A$ 是切割公式，
那么要么（在左前提是公理时）$!A \in \Gamma$ 且
$\Gamma = \Gamma', !A$，要么（在右前提是公理时）$!A \in
\Delta$ 且 $\Delta = \Delta', !A$。在第一种情况下，$\pi_2$ 是
$!A, !A, \Gamma' \Sequent \Delta$ 的一个推导，根据 \LeftR{\Contraction} 的可容许性（\olref[seq][inv]{prop:G3c-cont-adm}），存在
$\Gamma \Sequent \Delta$ 的一个推导。在第二种情况下，$\pi_2$ 是
$\Gamma \Sequent \Delta', !A, !A$ 的一个推导，根据 \RightR{\Contraction} 的可容许性，可得 $\Gamma \Sequent
\Delta$ 的一个推导。

\paragraph{情况 B：} 某个前提是形如~$\lfalse, \Gamma' \Sequent \Delta'$ 的公理。如果左前提是这样的
公理，那么 $\lfalse \in \Gamma$ 且 $\Gamma \Sequent \Delta$ 也是
公理。如果右前提是这样的公理，那么要么
$\lfalse \in \Gamma$（此时 $\Gamma \Sequent \Delta$ 仍是公理），要么 $!A \ident \lfalse$。在后一种情况下，$\pi_1$ 是
$\Gamma \Sequent \Delta, \lfalse$ 的一个推导。我们可以通过删除~$\pi_1$ 中每个相继式后件中的一个~$\lfalse$ 副本，将其转化为~$\Gamma \Sequent \Delta$ 的推导。
（练习：通过对~$\pheight{\pi_1}$ 归纳来验证这一点。）

现在假设情况~A 和情况~B 均不适用。在剩余的情况下，$\cheight{\pi} > 0$，即至少有一个前提是某个推理的结论。

\paragraph{情况 C 和 D：置换切割。}
情况~C 和~D 所考虑的情形都有一个共同模式。我们从一个以~\CutCS{} 结尾的推导开始，它的某个前提由规则~$R$ 给出，其中切割公式~$!A$ 不是主公式（某个其他公式~$!D$ 才是）。在这个推导中，
\CutCS{} 跟在规则~$R$ 之后，其示意图如下：
\[
\AxiomC{}
\RightLabel{$\pi_1$}
\Deduce$\Gamma \fCenter \Delta', !D, !A$
\AxiomC{}
\RightLabel{$\pi_3$}
\Deduce$!A, \Gamma, \Pi \fCenter \Delta', \Lambda$
\RightLabel{$R$}
\UnaryInf$!A, \Gamma \fCenter \Delta', !D$
\RightSubproofLabel{$\pi_2$}
\RightLabel{\CutCS}
\BinaryInf$\Gamma \fCenter \Delta', !D$
\DisplayProof
\]
这仅覆盖了 $R$ 为单前提右规则、$A$ 不是 $R$ 的主公式、而真正是主公式的 $!D$ 出现在~\CutCS{} 右前提后件中的情况。如果 $!D$~出现在前件中（即 $R$ 是左规则）或出现在~\CutCS 的左前提中，示意图当然会不同。
$\Pi$ 和~$\Lambda$ 中的公式代表规则~$R$ 的活动公式。

我们将这一顺序倒转，考虑一个 $R$ 跟在~\CutCS{} 之后的推导：
\[
\AxiomC{}
\RightLabel{$\pi_1'$}
\Deduce$\Gamma, \Pi \fCenter \Delta', !D, \Lambda, !A$
\AxiomC{}
\RightLabel{$\pi_3'$}
\Deduce$!A, \Gamma, \Pi \fCenter \Delta', !D, \Lambda$
\RightLabel{\CutCS}
\BinaryInf$\Gamma, \Pi \fCenter \Delta', !D, \Lambda$
\RightLabel{$R$}
\UnaryInf$\Gamma, \fCenter \Delta', !D, !D$
\DisplayProof
\]
我们分别从 $\pi_1$ 和 $\pi_3$ 出发，利用保持高度的弱化得到 $\pi_1'$ 和 $\pi_3'$。

由于我们交换了 \CutCS{} 与~$R$ 的顺序，这被称为“将切割向上穿过~$R$ 进行置换”。这样做降低了以 \CutCS{} 右前提为结束相继式的子推导的高度，从而降低了切割高度。也就是说，以新的 \CutCS{} 结尾的子推导的高度可假定不大于
$\pi_1$ 和 $\pi_3$，因此切割高度降低了（我们在右侧移除了一个推理）。现在归纳假设适用，由此可知这个 \CutCS{} 推理可以被移除。最后，我们利用~$\RightR{\Contraction}$ 的可容许性移除额外的~$!D$ 副本。

\paragraph{情况 C：}
$\pi_1$ 和 $\pi_2$ 中至少有一个以单前提规则~$R$ 结尾，且其中 $!A$~不是主公式。总共有 18 种情况，取决于 $!A$~在~\CutCS 的左前提还是右前提中不是主公式，以及规则~$R$ 是九个单前提规则（\LeftR{\lnot}, \RightR{\lnot}, \RightR{\lor}, \LeftR{\land},
\RightR{\lif}, 以及四个量词规则）中的哪一个。我们只考虑两个例子：$!A$~在~$\pi_2$ 的最后推理中不是主公式，且该推理以~$\RightR{\lif}$ 或~$\LeftR{\lexists}$ 结束。其他情况类似，留作练习。

假设 $\pi$ 以如下推理结尾：
\[
\AxiomC{}
\RightLabel{$\pi_1$}
\Deduce$\Gamma \fCenter \Delta', !B \lif !C, !A$
\AxiomC{}
\RightLabel{$\pi_3$}
\Deduce$!A, !B, \Gamma \fCenter \Delta', !C$
\RightLabel{$\RightR{\lif}$}
\UnaryInf$!A, \Gamma \fCenter \Delta', !B \lif !C$
\RightSubproofLabel{$\pi_2$}
\RightLabel{\CutCS}
\BinaryInf$\Gamma \fCenter \Delta', !B \lif !C$
\DisplayProof
\]
其中 $\Delta = \Delta', !B \lif !C$。我们对 $\pi_1$ 应用保持高度的弱化（\olref[seq][adm]{prop:weak-G3c-adm}），得到一个推导 $\pi_1'$，其结束相继式为 $!B, \Gamma \fCenter \Delta', !B \lif !C, !C,
!A$（我们在左侧弱化 $!B$，在右侧弱化 $!C$）。
类似地，我们对~$\pi_3$ 应用 \olref[seq][adm]{prop:weak-G3c-adm}，得到一个推导~$\pi_3'$，其结束相继式为 $!A, !B, \Gamma \fCenter
\Delta', !B \lif !C, !C$（通过在右侧弱化 $!B \lif !C$）。
归纳假设适用于以下推导（切割公式为~$!A$）：
\[
\AxiomC{}
\RightLabel{$\pi_1'$}
\Deduce$!B, \Gamma \fCenter \Delta', !B \lif !C, !C, !A$
\AxiomC{}
\RightLabel{$\pi_3'$}
\Deduce$!A, !B, \Gamma \fCenter \Delta', !B \lif !C, !C$
\RightLabel{\CutCS}
\BinaryInf$!B, \Gamma \fCenter \Delta', !B \lif !C, !C$
\DisplayProof
\]
该推导与~$\pi$ 具有相同的切割秩，但切割高度更低。由此可得一个在 \Log{G3c} 中不含 \CutCS{} 的推导~$\pi_4$，其结束相继式为~$!B, \Gamma \fCenter \Delta', !B \lif
!C, !C$。应用 $\RightR{\lif}$ 规则可得结束相继式为 $\Gamma \fCenter
\Delta', !B \lif !C, !B \lif !C$ 的推导：
\[
\AxiomC{}
\RightLabel{$\pi_4$}
\Deduce$!B, \Gamma \fCenter \Delta', !B \lif !C, !C$
\RightLabel{$\RightR{\lif}$}
\UnaryInf$\Gamma \fCenter \Delta', !B \lif !C, !B \lif !C$
\DisplayProof
\]
为得到 $\Gamma \Sequent \Delta$，即 $\Gamma
\fCenter \Delta', !B \lif !C$ 的推导~$\pi'$，需应用
\olref[seq][inv]{prop:G3c-cont-adm}（收缩规则的可容许性）。

现在假设 $\pi$ 以如下推理结尾：
\[
\AxiomC{}
\RightLabel{$\pi_1$}
\Deduce$\lexists[x][!B(x)], \Gamma' \fCenter \Delta, !A$
\AxiomC{}
\RightLabel{$\pi_3$}
\Deduce$!A, !B(c), \Gamma' \fCenter \Delta$
\RightLabel{$\LeftR{\lexists}$}
\UnaryInf$!A, \lexists[x][!B(x)], \Gamma' \fCenter \Delta$
\RightSubproofLabel{$\pi_2$}
\RightLabel{\CutCS}
\BinaryInf$\Gamma' \fCenter \Delta$
\DisplayProof
\]
我们再次将归纳假设应用于
\[
\AxiomC{}
\RightLabel{$\pi_1[!B(c) \Sequent {}]$}
\Deduce$\lexists[x][!B(x)], !B(c), \Gamma' \fCenter \Delta, !A$
\AxiomC{}
\LeftLabel{$\pi_3[\lexists[x][!B(x)] \Sequent {}]$}
\Deduce$!A, \lexists[x][!B(x)], !B(c), \Gamma' \fCenter \Delta$
\RightLabel{\CutCS}
\BinaryInf$\lexists[x][!B(x)], !B(c), \Gamma' \fCenter \Delta$
\RightLabel{\LeftR{\lexists}}
\UnaryInf$\lexists[x][!B(x)], \lexists[x][!B(x)], \Gamma' \fCenter \Delta$
\DisplayProof
\]
注意特征变元条件成立，因为~$\pi_2$ 中~\LeftR{\lexists} 的特征变元条件保证了
$c$ 不在 $\lexists[x][!B(x)]$、$\Gamma'$ 或~$\Delta$ 中出现。
最后，应用 \olref[seq][inv]{prop:G3c-cont-adm}。


\paragraph{情况 D：}
$\pi_1$ 和 $\pi_2$ 中至少有一个以双前提规则~$R$ 结尾，且其中 $!A$ 不是主公式。共有 6 种情况。我们考虑 $!A$ 在~$\pi_2$ 的最后推理中不是主公式，且该推理以~$\RightR{\land}$ 结束的情况；其他情况类似。

假设 $\pi$ 以如下推理结尾：
\[
\AxiomC{}
\RightLabel{$\pi_1$}
\Deduce$\Gamma \fCenter \Delta', !B \land !C, !A$
\AxiomC{}
\RightLabel{$\pi_3$}
\Deduce$!A, \Gamma \fCenter \Delta', !B$
\AxiomC{}
\RightLabel{$\pi_4$}
\Deduce$!A, \Gamma \fCenter \Delta', !C$
\RightLabel{$\RightR{\land}$}
\BinaryInf$!A, \Gamma \fCenter \Delta', !B \land !C$
\RightSubproofLabel{$\pi_2$}
\RightLabel{\CutCS}
\BinaryInf$\Gamma \fCenter \Delta$
\DisplayProof
\]
归纳假设适用于以下推导：
\[
\AxiomC{}
\RightLabel{$\pi_1'$}
\Deduce$\Gamma \fCenter \Delta', !B \land !C, !B, !A$
\AxiomC{}
\RightLabel{$\pi_3'$}
\Deduce$!A, \Gamma \fCenter \Delta', !B \land !C, !B$
\RightLabel{\CutCS}
\BinaryInf$\Gamma \fCenter \Delta', !B \land !C, !B$
\DisplayProof
\]
与
\[
\AxiomC{}
\RightLabel{$\pi_1''$}
\Deduce$\Gamma \fCenter \Delta, !B \land !C, !C, !A$
\AxiomC{}
\RightLabel{$\pi_4'$}
\Deduce$!A, \Gamma \fCenter \Delta', !B \land !C, !C$
\RightLabel{\CutCS}
\BinaryInf$\Gamma \fCenter \Delta', !B \land !C, !C$
\DisplayProof
\]
此处，\Log{G3c} 推导 $\pi_1'$ 和 $\pi_1''$ 是通过保持高度的弱化（\olref[seq][adm]{prop:weak-G3c-adm}）
从~$\pi_1$ 得到的（一次在右侧弱化 $!B$，一次弱化 $!C$）。推导 $\pi_3'$ 和~$\pi_4'$ 也是通过保持高度的弱化分别从 $\pi_3$ 和~$\pi_4$ 得到的（通过在右侧弱化 $!B \land !C$）。

归纳假设给出了 \Log{G3c} 中的推导~$\pi_5$ 和~$\pi_6$，其结束相继式分别为 $\Gamma \fCenter \Delta', !B \land !C, !B$ 和 $\Gamma
\fCenter \Delta', !B \land !C, !C$。我们使用
规则~$\RightR{\land}$ 将它们组合成结束相继式为 $\Gamma \Sequent \Delta', !B \land !C, !B \land !C$ 的推导：
\[
\AxiomC{}
\RightLabel{$\pi_5$}
\Deduce$\Gamma \fCenter \Delta', !B \land !C, !B$
\AxiomC{}
\RightLabel{$\pi_6$}
\Deduce$\Gamma \fCenter \Delta', !B \land !C, !C$
\RightLabel{\RightR{\land}}
\BinaryInf$\Gamma \fCenter \Delta, !B \land !C, !B \land !C$
\DisplayProof
\]
为得到 $\Gamma \Sequent \Delta$，即 $\Gamma \fCenter
\Delta', !B \land !C$ 的推导，需要应用 \olref[seq][inv]{prop:G3c-cont-adm}（收缩规则的可容许性）。

\paragraph{情况 E：归约切割。}
最后一种情况涉及 $\pi_1$ 和 $\pi_2$ 均以某个推理结尾，且在这些推理中 $!A$ 都是主公式。共有六种情况，对应于~$!A$ 的六种可能的主联结词。

在每种情况下，我们将一个切割公式为~$!A$ 的 \CutCS{} 转换成一个或多个对~$!A$ 的子公式（即那些以~$!A$ 为主公式的推理的活动公式）的切割。因此，这些情况通常被称为切割推理的“归约”或“转换”。

\begin{enumerate}
\item 若 $!A \ident \lnot !B$，则 $\pi$ 以如下推理结尾：
\[
\AxiomC{}
\RightLabel{$\pi_1'$}
\Deduce$!B, \Gamma \fCenter \Delta$
\RightLabel{\RightR{\lnot}}
\UnaryInf$\Gamma \fCenter \Delta, \lnot !B$
\LeftSubproofLabel{$\pi_1$}
\AxiomC{}
\RightLabel{$\pi_2'$}
\Deduce$\Gamma \fCenter \Delta, !B$
\RightLabel{\LeftR{\lnot}}
\UnaryInf$\lnot !B, \Gamma \fCenter \Delta$
\RightSubproofLabel{$\pi_2$}
\RightLabel{\CutCS}
\BinaryInf$\Gamma \fCenter \Delta$
\DisplayProof
\]
由归纳假设，可从以下推导中消除 \CutCS{}：
\[
\AxiomC{}
\RightLabel{$\pi_2'$}
\Deduce$\Gamma \fCenter \Delta, !B$
\AxiomC{}
\RightLabel{$\pi_1'$}
\Deduce$!B, \Gamma \fCenter \Delta$
\RightLabel{\CutCS}
\BinaryInf$\Gamma \fCenter \Delta$
\DisplayProof
\]

\item
若 $!A \ident (!B \lor !C)$，则 $\pi$ 以如下推理结尾：
\[
\AxiomC{}
\RightLabel{$\pi_1'$}
\Deduce$\Gamma \fCenter \Delta, !B, !C$
\RightLabel{\RightR{\lor}}
\UnaryInf$\Gamma \fCenter \Delta, !B \lor !C$
\LeftSubproofLabel{$\pi_1$}
\AxiomC{}
\RightLabel{$\pi_2'$}
\Deduce$!B, \Gamma \fCenter \Delta$
\AxiomC{}
\RightLabel{$\pi_2''$}
\Deduce$!C, \Gamma \fCenter \Delta$
\RightLabel{\LeftR{\lor}}
\BinaryInf$!B \lor !C, \Gamma \fCenter \Delta$
\RightSubproofLabel{$\pi_2$}
\RightLabel{\CutCS}
\BinaryInf$\Gamma \fCenter \Delta$
\DisplayProof
\]
考虑以 \CutCS 结尾的推导：
\[
\AxiomC{}
\RightLabel{$\pi_1'$}
\Deduce$\Gamma \fCenter \Delta, !B, !C$
\AxiomC{}
\RightLabel{$\pi_2'''$}
\Deduce$!C, \Gamma \fCenter \Delta, !B$
\RightLabel{\CutCS}
\BinaryInf$\Gamma \fCenter \Delta, !B$
\DisplayProof
\]
其中 $\pi_2'''$ 是通过保持高度的弱化从 $\pi_2''$ 得到的。由于 $\depth{!C} < \depth{!B \lor !C}$，其切割秩小于 $\cutr{\pi}$，
故由归纳假设可得 $\Gamma \fCenter \Delta, !B$ 的 \Log{G3c} 推导 $\pi_1''$。以 \CutCS 结尾的推导：
\[
\AxiomC{}
\RightLabel{$\pi_1''$}
\Deduce$\Gamma \fCenter \Delta, !B$
\AxiomC{}
\RightLabel{$\pi_2'$}
\Deduce$!B, \Gamma \fCenter \Delta$
\RightLabel{\CutCS}
\BinaryInf$\Gamma \fCenter \Delta$
\DisplayProof
\]
切割秩为 $= \depth{!B} < \depth{!B \lor !C}$，所以归纳假设再次适用，从而可得 $\Gamma
\fCenter \Delta$ 的推导 $\pi'$。

\item $!A \ident (!B \land !C)$。练习。

\item 若 $!A \ident (!B \lif !C)$，则 $\pi$ 以如下推理结尾：
\[
\AxiomC{}
\RightLabel{$\pi_1'$}
\Deduce$!B, \Gamma \fCenter \Delta, !C$
\RightLabel{\RightR{\lif}}
\UnaryInf$\Gamma \fCenter \Delta, !B \lif !C$
\LeftSubproofLabel{$\pi_1$}
\AxiomC{}
\RightLabel{$\pi_2'$}
\Deduce$\Gamma \fCenter \Delta, !B$
\AxiomC{}
\RightLabel{$\pi_2''$}
\Deduce$!C, \Gamma \fCenter \Delta$
\RightLabel{\LeftR{\lif}}
\BinaryInf$!B \lif !C, \Gamma \fCenter \Delta$
\RightSubproofLabel{$\pi_2$}
\RightLabel{\CutCS}
\BinaryInf$\Gamma \fCenter \Delta$
\DisplayProof
\]
以 \CutCS 结尾的推导：
\[
\AxiomC{}
\RightLabel{$\pi_1'$}
\Deduce$!B, \Gamma \fCenter \Delta, !C$
\AxiomC{}
\RightLabel{$\pi_2''$}
\Deduce$!C, \Gamma \fCenter \Delta$
\RightLabel{\CutCS}
\BinaryInf$!B, \Gamma \fCenter \Delta$
\DisplayProof
\]
切割秩低于~$\pi$。归纳假设给出
$!B, \Gamma \fCenter \Delta$ 的推导~$\pi_1''$。
现在考虑以 \CutCS 结尾的推导：
\[
\AxiomC{}
\RightLabel{$\pi_2'$}
\Deduce$\Gamma \fCenter \Delta, !B$
\AxiomC{}
\RightLabel{$\pi_1''$}
\Deduce$!B, \Gamma \fCenter \Delta$
\RightLabel{\CutCS}
\BinaryInf$\Gamma \fCenter \Delta$
\DisplayProof
\]
它的切割秩也小于~$\pi$ 的切割秩，所以归纳
假设给出 $\Gamma \fCenter
\Delta$ 的 \Log{G3c} 推导~$\pi'$。
\item 若 $!A \ident \lforall[x][!B(x)]$，则 $\pi$ 以如下推理结尾：
\[
\AxiomC{}
\RightLabel{$\pi_1'$}
\Deduce$\Gamma \fCenter \Delta, !B(c)$
\RightLabel{\RightR{\lforall}}
\UnaryInf$\Gamma \fCenter \Delta, \lforall[x][!B(x)]$
\LeftSubproofLabel{$\pi_1$}
\AxiomC{}
\RightLabel{$\pi_2'$}
\Deduce$!B(t), \lforall[x][!B(x)], \Gamma \fCenter \Delta$
\RightLabel{\LeftR{\lforall}}
\UnaryInf$\lforall[x][!B(x)], \Gamma \fCenter \Delta$
\RightSubproofLabel{$\pi_2$}
\RightLabel{\CutCS}
\BinaryInf$\Gamma \fCenter \Delta$
\DisplayProof
\]
由归纳假设，可从以下推导中消除 \CutCS{}：
\[
\AxiomC{}
\RightLabel{$\pi_1''$}
\Deduce$!B(t), \Gamma \fCenter \Delta, \lforall[x][!B(x)]$
\AxiomC{}
\RightLabel{$\pi_2'$}
\Deduce$!B(t), \lforall[x][!B(x)], \Gamma \fCenter \Delta$
\RightLabel{\CutCS}
\BinaryInf$!B(t), \Gamma \fCenter \Delta$
\DisplayProof
\]
其中 $\pi_1''$ 是通过保持高度的弱化从 $\pi_1$（而非~$\pi_1'$！）得到的。（该推导具有相同的切割秩但切割高度更低。）由此可得 $!B(t), \Gamma \fCenter \Delta$ 的推导~$\pi_2''$。

推导~$\pi_1'$ 的结束相继式为 $\Gamma \Sequent \Delta,
!B(c)$，其中 $c$~是某个~$\RightR{\lforall}$
推理的特征变元。所以 $c$~不在~$!B(x)$、$\Gamma$ 或~$\Delta$ 中出现。我们
假设推导~$\pi$ 是正则的，所以 $c$~仅是紧接其下的~$\RightR{\lforall}$ 推理的特征变元，并且只出现在该推理之上，即仅出现在~$\pi_1'$ 中，且不是~$\pi_1'$ 中任何其他推理的特征变元。由于 $c$ 不在~$\pi_2$ 中出现，它也不在~$t$ 中出现。根据 \olref[seq][qua]{lem:subst-regular}，
推导~$\Subst{\pi_1'}{t}{c}$ 是 $\Gamma \Sequent \Delta,
!B(t)$ 的一个推导。现在我们对以下推导应用归纳假设：
\[
\AxiomC{}
\RightLabel{$\Subst{\pi_1'}{t}{c}$}
\Deduce$\Gamma \fCenter \Delta, !B(t)$
\AxiomC{}
\RightLabel{$\pi_2''$}
\Deduce$!B(t), \Gamma \fCenter \Delta$
\RightLabel{\CutCS}
\BinaryInf$\Gamma \fCenter \Delta$
\DisplayProof
\]
（归纳假设适用，因为切割公式是~$!B(t)$，所以该推导的切割秩低于~$\pi$。）
\item 我们将 $!A \ident \lexists[x][!B(x)]$ 的情形留作练习。
\end{enumerate}
\end{proof}

\begin{prob}
验证 \olref{lem:cut-adm-G3c} 证明中情况~D 的一个子情况：在 $\pi_1$ 的最后一个推理（以 $\LeftR{\lif}$ 结束）中，公式 $!A$ 不是主公式。（注：本练习的部分要求是清楚地陈述你要证明什么，即在此情况下 $\pi$ 具有何种形式。）
\end{prob}

\begin{prob}
验证 \olref{lem:cut-adm-G3c} 证明中情况~D 的一个子情况：在 $\pi_1$ 的最后一个推理（以 $\LeftR{\lforall}$ 结束）中，公式 $!A$ 不是主公式。
\end{prob}

\begin{prob}
验证 \olref{lem:cut-adm-G3c} 证明中情况~E 的一个子情况：公式 $!A$ 在 $\pi_1$ 和 $\pi_2$ 各自的最后一个推理中都是主公式，且 $!A \ident (!B \land !C)$。
\end{prob}

\begin{prob}
验证 \olref{lem:cut-adm-G3c} 证明中情况~E 的一个子情况：公式 $!A$ 在 $\pi_1$ 和 $\pi_2$ 各自的最后一个推理中都是主公式，且 $!A \ident \lexists[x][!B(x)]$。
\end{prob}

注意，在情况~E 中，我们对以 \CutCS{} 结束的推导应用归纳假设，而该推导的切割高度可能大于原始证明的切割高度。例如，在 $!A \ident (!B \lif !C)$ 的情形下，我们将 $\pi$ 转变为一个包含两个 \CutCS{} 推理的推导：
\[
\AxiomC{}
\RightLabel{$\pi_2'$}
\Deduce$\Gamma \fCenter \Delta, !B$
\AxiomC{}
\RightLabel{$\pi_1'$}
\Deduce$!B, \Gamma \fCenter \Delta, !C$
\AxiomC{}
\RightLabel{$\pi_2''$}
\Deduce$!C, \Gamma \fCenter \Delta$
\RightLabel{\CutCS}
\BinaryInf$!B, \Gamma \fCenter \Delta$
\RightSubproofLabel{$\pi_1''$}
\RightLabel{\CutCS}
\BinaryInf$!B, \Gamma \fCenter \Delta$
\DisplayProof
\]
$\pi_1'$ 和 $\pi_2''$ 之间的顶层 \CutCS{} 的切割秩和切割高度均低于 $\pi$。但是，应用归纳假设的结果——推导 $\pi_1''$——的切割高度不再低于 $\pi$。然而，由于下层 \CutCS{} 的切割公式是 $!B$，其切割\emph{秩}低于 $\pi$ 的切割秩。

\olref{lem:cut-adm-G3c} 确立了我们可以从一个推导中消除单个最顶层的 \CutCS{} 推理。通过依次将 \olref{lem:cut-adm-G3c} 应用于以 \CutCS{} 结束的最顶层子推导，可以证明，我们能够从推导中消除任意数量的 \CutCS{} 推理。

\begin{thm}\ollabel{thm:cut-elim-G3c-topmost}
$\Log{G3c} + \CutCS$ 中的任意推导 $\pi$ 都可转换为 $\Log{G3c}$ 中的推导。
\end{thm}

\begin{proof}
对 $\pi$ 中 \CutCS{} 推理的数量 $n$ 进行归纳。若 $n=0$，则无需做任何事：$\pi$ 本身已是 $\Log{G3c}$ 中的推导。否则，考虑以最顶层 \CutCS{} 推理结束的子推导 $\pi'$。它包含单个 \CutCS{} 规则作为最后一个推理。应用 \olref{lem:cut-adm-G3c} 将其转换为无 \CutCS{} 的证明 $\pi''$，并用 $\pi''$ 替换子推导 $\pi'$。所得结果是一个包含 $n-1$ 个 \CutCS{} 推理的推导，从而可应用归纳假设。
\end{proof}

\end{document}